\section{Partition functions}\label{CLOSED}

Assume that $W$ is a {\it closed} oriented $3$-manifold, and that $L$
is a link in $W$ with a framed part $L_\Ff$ and an unframed one
$L^0$. Each variant of quantum hyperbolic field theory (see Section
\ref{variant}) leads to the respective {\it partition functions}. 

If $\rho$ is trivial at each meridian of $L^0$, we have the
QHFT {\it partition functions} 
$$ \Hh_N(W,L_\Ff,L^0,\rho,h,k)$$
that specialize to the QHFT$^0$ ones when $L=L^0$:
$$H_N(W,L,\rho,h,k)= \Hh_N(W,\emptyset,L^0,h,k) \ .$$
If $\rho$ is also assumed to be non trivial at each meridian of
$L_\Ff$, we have also
$$\Hh_N^e(W,L_\Ff,L^0,\rho,h,k) \ .$$ These partition functions are
scalars, well-defined up to sign and multiplication by $N$th-roots of
unity. Typical examples of triples $(W,L_\Ff,\rho)$ are given by
hyperbolic {\it cone} manifolds $W$ with framed cone locus $L_\Ff$ and
hyperbolic holonomy $\rho$ on $W \setminus L_\Ff$. The partition
functions can be expressed in terms of manifolds $Y$ with toric
boundary and containing an unframed link $L^0$ in the interior. By
fixing an ordered basis $(m_i,l_i)$ for the integral homology of each
boundary torus, let $W$ be obtained from $Y$ by Dehn filling along the
$m_i$, and $L_\Ff$ be the disjoint union of the cores of the filling
solid tori, framed by the $l_i$. Then the partition functions of
$(W,L_\Ff,L^0,\rho,h,k)$ are in fact invariants of
$(Y,\{(m_i,l_i)\}_i,L^0,\rho,h,k)$.

\subsection{QHFT vs QHFT$^0$ partition functions}\label{ALL}
For $\Bb=(W,L,\rho)$, $L=L^0$, with weights $h=k=0$,
$H_N(W,L,\rho,0,0)$ coincide with the invariants $H_N(W,L,\rho)$
constructed in \cite{BB1, BB2}.  Let us consider more generally
$(W,L_\Ff,L^0,\rho,0,0)$.  Fix also a framing $\Ff_0$ for
$L^0$. Then we can consider the partition function
$\Hh_N(W,L_\Ff\cup L^0_{\Ff_0},\emptyset,\rho,0,0)$.  Let us denote by
$\bar \lambda$ the unframed link obtained by splitting each component
of $L^0$ in the two corresponding parallel boundary components of the
ribbon link $L^0_{\Ff_0}$. We have:
\begin{prop}\label{QHFTQHI} 
$$ \Hh_N(W,L_\Ff\cup L^0_{\Ff_0},\emptyset,\rho,0,0) =_N
\Hh_N(W,L_\Ff, \bar \lambda,\rho,0,0) \ . $$ 
\end{prop}

\noindent {\it Proof.} For simplicity, assume that $L=L^0$.  Fix a
$\Dd$-triangulation of $(\tilde{W},\tilde{L}_\Ff,\rho)$ where each
tunnel component $B$ has a symmetric admissible triangulation as in
Figure \ref{carre} (opposite sides of the quadrilateral are
identified). The tangle $\bar{\lambda}$ cuts open $B$ into symmetric
annuli, left and right to the central vertical line in Figure
\ref{carre}.
\begin{figure}[ht]
\begin{center}
\includegraphics[width=3cm]{carre.eps}
\caption{\label{carre} A special admissible triangulation of $B$.}
\end{center}
\end{figure}

Because $\rho$ has trivial holonomy at the meridians of $L_\Ff$,
we can assume that the cocycle takes the same values on symmetric
edges. Identifying the annuli we thus get a $\Dd$-triangulation
for the QHFT$^0$ triple $(W,\bar{\lambda},\rho)$. Since
$H_N(W,\bar{\lambda},\rho)$ is computed from the idealization and
symmetric tetrahedra in the cusps have opposite branching
orientation, the result will follow if we show the existence of
symmetric flat/charges. Then each cusp tensor will be the identity
map.

The existence of flattenings with this property can be shown using
Remark \ref{canonicalflat}, but for charges we need to take
another route. Recall from the end of section \ref{QHG} that
flat/charges form affine spaces over an integral lattice generated
by vectors attached to the edges. For an edge of $B$, such vectors
can be represented as adding $+1$ at one of the left adjacent
corner and $-1$ at the other, and the inverse for the right
adjacent corners. Using these rules and (\ref{2-3comp}) it is
straightforward (though tedious) to check that any given
flat/charge can be turned into one with equal quantum log-branches
on symmetric tetrahedra.\hfill$\Box$

\subsection{Invariants of cusped manifolds and surgery formulas}\label{DEHNF}

Let us recall the QHG pseudo-manifold triangulations $\Tt$
used in \cite{BB2} (see Definition 6.2 and Definition 6.3 in that
paper) to define the quantum hyperbolic invariant trace tensors
$\Hh_N(\Tt)$ for oriented cusped hyperbolic manifolds.\smallskip

Let $M$ be a cusped manifold. Denote $Z$ the pseudo-manifold obtained
by taking the one point compactification of each cusp of $M$. $M$
admits a triangulation by positively embedded hyperbolic ideal
tetrahedra, possibly including some degenerate ones of null volume
(ie. having real cross-ratios). Such a triangulation can be obtained
by subdividing the canonical Epstein-Penner cell decomposition of
$M$. This gives rise to triangulation $(T_0,z_0)$ of $Z$, where $z_0$ is
the cross-ratio function of the abstract edges of $T$, the imaginary
part of every cross-ratio being $\geq 0$. We call it a {\it
  quasi-geometric} ideal triangulation of $Z$. If some quasi-geometric
triangulation admits a global branching, we say that $M$ is {\it
  gentle}. More generally, $M$ is said {\it weakly-gentle} if there is
an $\Ii$-triangulation $(T,b,w)$ of $Z$ such that $(T,z)$,
$z=w^{*_b}$, is obtained via a (possibly empty) finite sequence
$(T_0,z_0) \rightarrow \ldots \rightarrow (T_i,z_i) \rightarrow \ldots
\rightarrow (T,z)$ of positive $2 \rightarrow 3$ transits, where
$(T_0,z_0)$ is a quasi-geometric triangulation of $Z$ as above. Each
transit $(T_i,z_i) \rightarrow (T_{i+1},z_{i+1})$ is defined by
$W_{T_i}(e)=W_{T_{i+1}}(e)$, with all exponents $*_b=1$ (see
(\ref{totw})). Every such a $(T,b,w)$ can be enhanced to flat/charged
$\Ii$-triangulations $\Tt=(T,b,w,f,c)$.
\smallskip

Given $(T_0,z_0)$, it is certainly possible to get an
$\Ii$-triangulation $(T,b,w)$ by performing {\it also} some bubble
moves (hence introducing new interior vertices). The authors do not
know any example of non weakly-gentle cusped manifold, that is, such
that we are forced to do it. Anyway, dealing with bubble moves is a technical
difficulty which will appear also in the proof of Theorem
\ref{defscissors} (1) below. We overcome it as follows. We fix an edge
$a$ of the canonical Epstein-Penner cell decomposition of $M$, and
take $A$ made by two copies of $a$ that intersect at non manifold
points of $Z$; the second copy runs parallel to $a$ within an open
cell of the decomposition.  Hence, $A$ is a circle covered by two
arcs.  We need to enlarge the notions introduced in Definition
\ref{DISFLATCHARGE}. We say that {\it $(T,H)$ is a distinguished
  triangulation of $(Z,a)$} if $H$ is a subcomplex of the $1$-skeleton
of $T$ isotopic to $A$, that contains all the regular vertices of $T$,
and such that one arc of $A$ is covered by an edge $l$ of $H$. We say
that {\it $c$ is a global charge on $(T,H)$} if
\begin{equation}\label{condc2}C_{\Tt}(e)=\left
\lbrace\begin{array}{l} 4\\0\\ 2\end{array}\right.\begin{array}{l}
{\rm if}\ e=l\\{\rm if}\ e \subset H\setminus l\\ {\rm if}\ e \subset
T \setminus H \end{array}
\end{equation}
If $H=\emptyset$ this reduces to the usual notion of global charge on
a closed triangulated pseudo manifold whose non manifold points have
toric links. By using bubble moves and the existence of such usual
global charges, it is easily seen that $(T,H)$ supports global charges
as in (\ref{condc2}) (see the proof of Theorem 6.8 in \cite{BB2} for
the details).

We say that $\Tt=(T,H,b,w,f,c)$ is a flat/charged $\Ii$-triangulation
of $(Z,a)$ if $(T,H,b,c)$ is a branched, charged and distinguished
triangulation of $(Z,a)$, and $(T,z)$, $z=w^{*_b}$, is obtained from a
quasi geometric $(T_0,z_0)$ via a finite sequence $(T_0,z_0)
\rightarrow \ldots \rightarrow (T,z)$ of transits supported by
positive $2 \rightarrow 3$ moves {\it and bubble moves}. By setting $a
= \emptyset$ and $H=\emptyset$, this definition incorporates that for
the weakly-gentle case.
\smallskip

In \cite{BB2} it is shown that flat/charged $\Ii$-triangulations $\Tt$
of $(Z,a)$ (with arbitrary weights) do exist and that, for every odd
$N\geq 1$,
$$ H_N(M,a)=\Hh_N(\Tt)$$ is a well defined invariant of $(M,a)$,
providing the weights of flat/charges to be $0$. To simplify the
exposition, below we continue with this normalization.  When $M$ is
weakly-gentle we get invariants $H_N(M)$. In fact, as a by product of
the following discussion, we will realize that $H_N(M,a)$ does not
depend on the choice of $a$, so that $H_N(M)$ is always well defined
(see Corollary \ref{DF-converge}).

\medskip

Let us recall now a few facts related to {\it hyperbolic Dehn filling}
(see \cite{Th}, \cite{BP}, \cite{Pe-Po}). A quasi geometric
triangulation $(T_0,z_0)$ as above corresponds to the complete
structure of $M$. It can be deformed in a complex variety of dimension
equal to the number of cusps. If $z'$ is close enough to $z_0$,
$(T_0,z')$ is a triangulation by (possibly negative - see
\cite{Pe-Po}) embedded hyperbolic ideal tetrahedra in a non-complete
hyperbolic structure, say $M'$, close to $M$.  In some case the
completion of $M'$ gives rise to a compact closed hyperbolic manifold
$W$, topologically obtained by Dehn filling of the (truncated) cusps
of $M$. The core of each attached solid torus is a ``short'' simple
closed geodesic $L_j$ of $W$, so that we have the (geodesic) link
$\textstyle L=\coprod_j L_j$. Moreover, there are sequences
$(W^n,L^n)$ obtained in this way such that the length of $L^n$ goes to
$0$ when $n\to +\infty$. Hence $(W^n,L^n)$ converges to the cusped
manifold $M$ (in a neat geometric sense). From now on we will consider
small deformations $z'$ leading to such closed completions.
\smallskip

As well as $(T_0,z_0)$ gives rise to a triangulation $\Tt=(T,H,b,w,f,c)$
of $(Z,a)$, $z'$ close to $z_0$ gives rise to another
flat/charged $\Ii$-triangulation $\Tt'=(T,H,b,w',f',c)$, where $w'$ is
close to $w$ and the log-branch associated to the global flattening
$f'$ corresponds to a continuous deformation of the one for $f$.

\begin{lem}\label{deformedlogb}{\rm (See \cite{N}, p. 469) Let $z'$ be
a small deformation of $z_0$ producing $(W,L)$, and $m_j$ be a
meridian of each link component $L_j$. Then there exist flattenings
$f''$ for the deformed triangulation $(T,b,w')$ such that the weight
$\gamma(f'')$ associated to the collection of log-branches of
$\Tt''=(T,H,b,w',f'',c)$ satisfies $\gamma(f'')(m_j)=0$ for all $j$.}
\end{lem}
\noindent {\it Proof.} The logarithm of the derivative of the
holonomy of $m_j$ is $0$ at the complete structure, but after
deformation it represents a full
$2\pi$-rotation about $L_j$ (see eg. \cite{Th} or \cite{BP}).
Hence $\gamma(f')(m_j) \in \mz\sqrt{-1}\pi$. The result then follows from
the comments after Theorem \ref{theoweight}.\hfill$\Box$
\smallskip

An explicit construction of flattenings $f''$ as in the lemma shall be
recalled during the proof of theorem \ref{defscissors}. Hence, we
dispose of two flat/charged $\Ii$-triangulations of $(Z,a)$, $\Tt'=
(T,H,b,w',f',c)$ and $\Tt''=(T,H,b,w',f'',c)$, relative to a small
deformation $z'$ of the complete structure $z_0$ as above, leading to
respective trace tensors $\Hh_N(\Tt')$ and $\Hh_N(\Tt'')$. We can now
state the main results of this Section. 

\begin{teo}\label{defscissors}
{\rm[Cusped manifold surgery formula]} Let $(W,L)$ be obtained by
completion of a small deformation $z'$ of $z_0$, and $\Tt'$, $\Tt''$
be associated triangulations. Denote by $\rho$ the hyperbolic holonomy
of $W$. Then we have $H_N(W,L,\rho) =_N
\Hh_N(\Tt'')$. Moreover, associated to each cusp $C_j$ of $M$ there is
an explicitely known map $\Lambda_N^j(\Tt''): \{ N-{\rm states\ of\ T}
\}\rightarrow \mc$ such that the following {\rm surgery formula}
holds:
$$H_N(W,L,\rho) =_N \sum_s \prod_{\Delta \subset T}
    \Rr_N(\Delta,b,w',f',c)_s\ \prod_j\Lambda_N^j(\Tt'')(s)$$ where
    $s$ runs over the $N$-states of $T$ and
    $\Rr_N(\Delta,b,w',f',c)_s$ is the matrix dilogarithm entry
    determined by $s$, for the tetrahedron with the continuously
    deformed structure.
\end{teo}
\begin{cor}\label{DF-converge} 
If $\{(W_n,L_n,\rho_n)\}$ is a sequence of closed hyperbolic Dehn
fillings converging to the cusp manifold $M$, then for every arc $a$ we
have $\textstyle \lim_n H_N (W_n,L_n,\rho_n)=_N H_N(M,a)$. Hence
$H_N(M)=H_N(M,a)$ is always a well defined invariant of $M$ (beyond
the weakly-gentle case).
\end{cor}

\begin{remarks}\label{CS-surg}
{\rm (1) Theorem \ref{defscissors} is the analog for $N>1$ of Theorems
  14.7 and Theorem 14.5 in \cite{N}, which describe surgery formulas
  for the volume, Vol($W$), and Chern-simons invariant, CS($W$), of
  $W$:
$$\sqrt{-1}({\rm Vol}(W) + \sqrt{-1} {\rm CS}(W)) = \sum_{\Delta
  \subset T} \mathcal{R}(\Delta,b,w',f')-\frac{\pi\sqrt{-1}}{2}
\sum_j \lambda(L_j)$$ where $\Rr$ is given by (\ref{symRmatdil})
and $\lambda(L_j)$ is the complex length of $L_j$, that is, the logarithm of
the dilation factor of its holonomy, which is a loxodromic
transformation of $\mh^3$. The technical complications due to the bubble
moves disappear for $N=1$.
\smallskip

(2) If $M$ is gentle and has a geometric branched ideal
triangulation $(T,b,w)$ without degenerate tetrahedra, then for each
$3$-simplex the flattenings of $\Tt'$ for a sufficiently small
deformation are just $-*_b$ times integral charges. It follows from
the proof of Theorem \ref{defscissors} that the scissors congruence
class $\mathfrak{c}_{\Ii_{fc}}(W,L,\rho)$ of \cite{BB2}, section 7,
coincides with Neumann's deformed scissors congruence class
$\hat{\beta}(M')$ in \cite{N}, Theorem 14.7 (see also Remark 6.12 and
Conjecture 7.9 in \cite{BB2}, where the undeformed $\hat{\beta}(M)$ is
denoted $\mathfrak{c}_{\Ii_{fc}}(M)$).
\smallskip

(3) In general there are small deformations $z$ of $z_0$ leading to
complete manifolds that are still cusped, that is only some cusps of
$M$ have supported a hyperbolic Dehn filling. There are also sequences
of such cusped manifolds $M^n$, with (short) geodesic links $L^n$,
converging to $M$. Similarly to the fusion of QHFT with QHFT$^0$
(see Section \ref{variant}) we can define quantum hyperbolic
invariants $H_N(M^n,L^n)$ for which the natural extensions of Theorem
\ref{defscissors} and Corollary \ref{DF-converge} hold.}
\end{remarks}
Let us consider now $(W,L_\Ff,L^0,\rho)$. Let $L_j$ be a component
of $L_\Ff$, $\bar \lambda_j = L_j \cup L_j'$, $L'_j$ being
the longitude of $L_j$ specifying the framing. Let $U=U(L_j)$
be a tubular neighbourhooh of $L_j$ in $W$, and $l\subset \partial U$
be a non separating simple closed curve. Let $W(l)$ be obtained from
$W$ by the Dehn filling of $W\setminus {\rm Int}(U)$ along $l$. Denote by $l^*$
the core of the attached solid torus.
\begin{teo} \label{surcomp} {\rm
[Closed manifold surgery formula]} Assume that $\rho(l)={\rm id} \in
  PSL(2,\mc)$ and the weight $k$ satisfies $k([l])=0$. Denote: $\rho'$
  the natural extension of $\rho_{|W \setminus U}$ to $W(l)$;
  $\widetilde L_\Ff = L_\Ff \setminus L_j$; $k'$ the restriction of
  $k$ to $W(l)$. We have
\begin{equation}\label{eqsurg1}
\Hh_N(W,L_\Ff,L^0,\rho,0,k) =_N \ \Hh_N(W(l),\widetilde L_\Ff, L^0\cup l^* \cup
\bar{\lambda}_j,\rho',0,k') \ .
\end{equation}
If moreover $\rho$ is not trivial at the meridians of $L_\Ff$ and $l$
is a longitude of $\partial U$, then

\begin{equation}\label{eqsurg2}
\Hh_N^e(W,L_\Ff,L^0,\rho,0,k) =_N \Hh_N^e(W(l),\widetilde L_\Ff, L^0\cup l^*
\cup L'_j,\rho',0,k') \ .
\end{equation}
\end{teo}
Let us assume now that $L^0$ is made by $r$ parallel copies of $L_j$
along the ribbon $L_\Ff$ that encodes the framing. So denote $L^0$ by
$\lambda_r$; with this notation, $\bar \lambda = \lambda_2$. Assume furthemore
that $l=m$ is a meridian of $L_j$, so that $l^*=L_j$. By appling inductively
both (\ref{eqsurg1}) and Proposition \ref{QHFTQHI} to this situation we get
\begin{cor} For every $r\geq 2$ we have
$$H_N(W,\lambda_r,\rho)=_N H_N(W,\lambda_2,\rho) \ . $$
\end{cor}


\begin{remarks}\label{remsurg}{\rm  
\smallskip

(1) Though disjoint and complementary by hypothesis, formula
    (\ref{eqsurg2}) is formally the same as that of Proposition
    \ref{QHFTQHI}, when replacing $l$ by $m$.
\smallskip

(2) Assume (for simplicity) that $L=L_\Ff$. When $l$ is a longitude of
$L_j$, $l^*$ inherits a natural framing in $W(l)$. Hence we get a
triple $(W(l),\hat L_{\hat \Ff},\rho)$. It follows from the very
definition of the QHFT tensors that $\Hh_N(W,L_\Ff,\rho,0,0)=
\Hh_N(W(l),\hat L_{\hat \Ff},\rho,0,0)$ and the same with $\Hh_N^e$
(when defined) replacing $\Hh_N$.

\smallskip

(3) We have seen that both $\Hh_N$ and $\Hh_N^e$ partitions functions
display interesting features of QHG. A main advantage of the $\Hh_N$
ones is the possibility to set in a {\it same ``holomorphic family''}
the QHG tensors associated to characters that are both trivial and non
trivial at link meridians.  Consider for example a hyperbolic knot $L$
in $S^3$, endowed with the canonical framing $\Ff$. Kashaev's volume
conjecture concerns the asymptotic behaviour of $H_N(S^3,L,\rho_{{\rm
    triv}})$ when $N\to +\infty$.  A reasonable variant of it is in
terms of the partition functions $\Hh_N(S^3,L_\Ff,\rho_{{\rm triv}}
,0,0)=_N H_N(S^3,\bar \lambda,\rho_{{\rm triv}}) $.  A family as above
could be useful in order to establish connections with the $\Hh_N$
partition functions of $(S^3,L_\Ff,\rho_{{\rm hyp}})$, where
$\rho_{\rm hyp}$ is the hyperbolic holonomy of the cusped manifold $M=
S^3\setminus L$.}
\end{remarks}

The rest of the section is devoted to the proof of these results.
This goes in several steps.

\smallskip

A main tool is the {\it simplicial blowing up/down} procedure
considered by Neumann in \cite{N}, section 11. We use it just
to get a simplicial version of (topological) Dehn filling.  Let $Z$ be
a pseudo manifold without boundary such that every non-manifold point
has toric link. Let $v$ be a non-manifold point. Consider a closed
cone neighborhood $N(v)$ of $v$, and a non separating simple closed
curve $C$ on the torus  $\partial N(v)$. The topological Dehn
filling of $Z$ at $v$ along $C$ is the pseudo manifold $Z'$ obtained by
gluing a $2$-handle to $Z\setminus {\rm Int}(N(v))$ along $C$, and then
collapsing to one point the resulting boundary component. 
\smallskip

Now, let $T$ be a pseudo manifold triangulation of $Z$. Consider the
{\it abstract star} Star$^0$($v$) of $v$ in $T$. The boundary of
Star$^0$($v$) is the {\it abstract} link Link($v$) which is
homeomorphic to $\partial N(v)$. Assume that the curve $C$ is realized
as a {\it simplicial} curve on Link($v$). Then the cone from $v$ over
$C$ in Star$^0$($v$) is a triangulated disk $D^0$. The interior of
Star$^0$($v$) embeds onto the interior of the actual {\it star} of $v$
in $T$, Star($v$), which is made of the union of the $3$-simplices
having $v$ as a vertex. In this way $D^0$ maps onto a triangulated
singular disk $D$ in $Z$, that has embedded interior and singular
boundary immersed in the boundary of Star($v$). Cut open $T$ along
${\rm Int}(D)$ and glue the double cone $CD$ of $D$ (this is a
triangulated singular $3$-ball, see Figure \ref{doublecone}) so that
the top and the bottom get identified with the two copies of $D$
resulting from slicing. This gives a triangulation $T'$ of the
pseudo-manifold $Z'$ obtained by Dehn filling along $C$. It has the
property that every (abstract) tetrahedron of $T$ persists in
$T'$. Referring to the topological description, the interior of the
co-core of the $2$-handle attached to $Z\setminus {\rm Int}(N(v))$ is
isotopic to the interior of the union $H'$ of two edges, each joining
$v$ to the new vertex $v'$ at the ``center" of $CD$. In fact $H'$ is
the core of the solid torus added by the Dehn filling.

\begin{remark}\label{fewsimp}{\rm 
Note that in general there are very few simplicial curves on a given
Link($v$). Hence to get such a simplicial description of an arbitrary
Dehn filling, we will usually have to modify a given
triangulation. For the peculiar QHG pseudo-manifold triangulations
considered in this section, retriangulating will be possible by using
QHG isomorphisms, hence without altering the trace tensors. In fact,
any two triangulations of $\partial N(v)$ are connected by a finite
sequence of $2$-dimensional $1 \leftrightarrow 1$ ``flip'' moves (see
Figure \ref{eTflip}), and $1 \leftrightarrow 3$ moves obtained by
replacing a $2$-simplex with the cone of its boundary to a
point. Since $N(v)$ is homeomorphic to Link($v$), any such a sequence
is the boundary trace of a sequence of $2\leftrightarrow 3$ moves and
bubble moves in Star($v$). (Note, in particular, that $H$ passes
through the new interior vertices). By using the arguments of Theorem
6.8 in \cite{BB2}, we will always be able to choose that sequence so
that it lifts to a sequence of QHG transits.}
\end{remark}
\begin{figure}[ht]
\begin{center}
\includegraphics[width=6cm]{doublecone.eps}
\caption{\label{doublecone} An instance of double cone on a disc.}
\end{center}
\end{figure}

Let us consider now a distinguished triangulation $(T,H)$ of $(Z,a)$
as above, and let $Z'$, triangulated by $T'$, be the result of a
simplicial Dehn filling of $Z$ along a curve $C$. Denote $H''$ the graph
union of the knot $H'$ (the core of the solid torus) and the image of
$H$ in $T'$. We define the notion of global charge on $(T',H'')$
by formally replacing $H$ by $H''$ in (\ref{condc2}).

\begin{lem}\label{extc} Let $c$ be a global charge on $(T,H)$
such that the charge weight of the curve $C$ is $0$. Then $c$ extends
to a global charge $c'$ on $(T',H')$.
\end{lem}
\noindent {\it Proof.} The complex $CD$ is made of pairs of
adjacent $3$-simplices, respectively above and below the disk $D$.
For a $3$-simplex of the top layer with charges $c_0$ and $c_1$ at
the edges in $D$ (ordered by using an orientation of $D$, say), we
will put the charges $-c_0$ and $2-c_1$ at these edges for the
symmetric $3$-simplex in the bottom layer. Then the other charges
are $c_2=1-c_0-c_1$ and $-c_2$, respectively. We have the charge
sum $C_{T'}(e)=2$ at each interior edge of $D$, and $C_{T'}(e) =
C_{T}(e)$ at the edges $e$ of $\partial D$. For the top edges $e'$
of $CD$ we can also choose the charges so that $C_{T'}(e')$ equals
$C_T(e)$, where $e$ is the copy of $e'$ in $D \subset M$. Indeed,
there are $n$ degrees of freedom in doing this, where $n$ is the
number of $1$-simplices in the curve $C$ used for blowing down. Then
we check that $C_{T'}(e')=2$ at the bottom edges. In particular,
the subcomplex $H$ survives in $T'$.

Note that $C_{T'}(e_0)=-C_{T'}(e_1)$ at the edges $e_0$ and $e_1$
of $H'$. We have to check that $C_{T'}(e_0)=0$, so that
(\ref{condc}) is satisfied on $(T',H''\setminus l)$. In fact,
$C_{T'}(e_0)$ is $n$ minus the sum of the $2n$ charges at the
bottom edges of $CD$, which is also the sum of charges in
$T'\setminus CD$ at these edges, minus $n$. We can form $n$ pairs
of such charges corresponding to the $3$-simplices of the ideal
triangulation $T$ of $M$ having a $2$-simplex in $D$. Replacing
for each of them the pair with $1$ minus the last charge, we get
that $C_{T'}(e_0)$ is equal to $\gamma(a)$, with $\gamma$ defined
in section \ref{QHG} and $a$ is a normal path in Link($v$) that
runs parallel to $C$ on one side (see Figure \ref{normalpath}).
Because the weight of $C$ is zero, we deduce $C_{T'}(e_0)=0$.

If $H \ne \emptyset$, we have to show that it can be deleted from
$H''$. As the two components $l$ and $H \setminus l$ are isotopic and
satisfy $C_{T'}(l)=4$ and $C_{T'}(e)=0$ for each edge $e \in H
\setminus l$, we can retriangulate the surgered pseudo-manifold $Z'$
so as to delete them, by using a sequence of charge transits starting
from $(T',H'')$ and terminating with a negative bubble move (see
\cite{BB1}, Proposition 4.27 and \cite{BB2}, proof of Theorem
6.8). Retriangulating $Z'$ backward, we eventually find a sequence of
charge transits terminating at $(T',H',c')$. The result
follows. \hfill$\Box$
\medskip
\begin{figure}[ht]
\begin{center}
\includegraphics[width=6cm]{normalpath.eps}
\caption{\label{normalpath} A normal path $a$ in Link($v$) running
parallel to a blowing-up curve $C$. Four $3$-simplices of $T'
\setminus CD$ glued to three $2$-simplices of $D$ are shown in the
picture.}
\end{center}
\end{figure}

\noindent {\it Proof of Theorem \ref{defscissors}.}  To simplify
assume that $M$ has only one cusp. Take $\Tt''=(T,H,b,w',f'',c)$. For
the first claim we assume that the meridian $m$ of $L$ is a simplicial
path in Link($v$), where the vertex $v$ of $T$ corresponds to the
filled cusp. This is possible due to Remark \ref{fewsimp}.

Lemma \ref{extc} implies that $c'$ extends to $(T',H')$ after the Dehn
filling along $m$. Extend the branching $b$ by letting the new vertex
$v'$ be a pit of the double cone $CD$ we splice in $T$. By using Lemma
\ref{deformedlogb}, arguments similar to that of Lemma \ref{extc} show
that we can give the same log-branches on the $3$-simplices of $CD$,
in a pair above and below the disk $D$ (see \cite{N}, p. 454). Hence
we get a distinguished flat/charged $\Ii$-triangulation for
$(W,L,\rho)$. The weight $h \in H^1(W;\mz/2\mz)$ is clearly $0$
because of the epimorphism $ H_1(M;\mz/2\mz) \rightarrow
H_1(W;\mz/2\mz)$ induced by inclusion.  As in the first claim of Lemma
\ref{CYLtensor} (2) we see that the (unnormalized) trace tensor for
$CD$ is $N$ times the identity map from the linear space attached to
the top copy of $D$ to that for the bottom one. Combining this with
the normalization of trace tensors in (\ref{deftt}) gives
$H_N(W,L,\rho) = \Hh_N(\Tt'')$.

By Lemma \ref{deformedlogb} we know that $\gamma(f'')-\gamma(f') \in
H^1(\partial M;\mz\pi\sqrt{-1})$ is non zero only at the class of $m$,
where it is $-2\sqrt{-1}\pi$. Hence the collection of values of
$f''-f'$ determines a path $l$ normal to the cusp triangulation
induced by $T$, that intersects $m$ once and whose homology class is
Poincar\'e dual to $(\gamma(f'')-\gamma(f'))/2\sqrt{-1}\pi$. Denote
$$\Delta(l_j)=(\Delta^1,\ldots,\Delta^{\mid \Delta(l_j) \mid})$$ the
sequence of flat/charged $\Ii$-tetrahedra (possibly with repetitions)
determined by the $2$-simplices met by $l$. Each time $l$ goes through
a $2$-simplex it selects one of its vertices, whence a cross-ratio
modulus, say $z_i$, of the tetrahedron $\Delta^i$ corresponding to the
$2$-simplex. The values of $f''$ on $\Delta^i$ are obtained from those
of $f'$ by adding or substracting $1$ at the edges corresponding to
the other two vertices, as indicated in Fig. \ref{surgcurve}. For any
fixed tetrahedron $\Delta$ of $T$ all three flattenings may be
eventually altered, and/or differ from those of $f'$ by adding or
substrating $n_i \in \mz$ with $n_i\ne-1$ or $1$, exactly when
$\Delta=\Delta^i=\Delta^j$ for some $i\ne j$. Now, recall from
(\ref{deftt}) that
$$H_N(M,a) = \Hh_N(\Tt') = \sum_s \prod_{\Delta \subset T}
\Rr_N(\Delta,b,w,f,c)_s.$$ Put $\zeta=\exp(2\pi\sqrt{-1}/N)$. For any
$x\in \mc \setminus \{\zeta^j,j=1,\ldots,N-1\}$ the function $g$
defined in the Appendix satisfies (\cite{BB2}, Lemma
8.2): $$g(x\zeta^k)=g(x)\prod_{j=1}^k
\frac{(1-x^N)^{1/N}}{1-x\zeta^j}.$$ Using this formula, it is easily
checked that given a flat/charged $\Ii$-tetrahedron
$(\Delta,b,w',f',c)$ with $f'=(f_0',f_1',f_2')$ and positive branching
orientation, if $f''=(f_0'+n,f_1',f_2'-n)$ we have
\begin{equation}\label{defqR} \Rr_N(\Delta,b,w',f'',c)_s
=\Rr_N(\Delta,b,w',f',c)_s\
\prod_{j=1}^n\frac{(w_1')^{-1}}{1-w_0'\zeta^{i-k+j}}
\end{equation}
up to multiplication by $N$th roots of unity, where $i$ and $k$ are as
in (\ref{symqmatdil}). For each $2$-simplex met by $l$ we can apply
(\ref{defqR}) to the corresponding tetrahedron, or the similar formula
(deduced from Corollary 5.6 of \cite{BB2}) for any other
branching. This defines the function $\Lambda_N(\Tt'')$, so that we get
$$\begin{array}{lll} \Hh_N(\Tt'') & = &
  \sum_s \prod_{\Delta \subset T} \Rr_N(\Delta,b,w',f'',c)_s\\ & = &
  \sum_s \prod_{\Delta \subset T}
  \Rr_N(\Delta,b,w',f',c)_s\ \Lambda_N(\Tt'').
\end{array}$$
The conclusion follows from the equality $H_N(W,L,\rho) =
\Hh_N(\Tt'')$.\hfill$\Box$
\begin{figure}[ht]
\begin{center}
\includegraphics[width=6cm]{surgcurve.eps}
\caption{\label{surgcurve} Flat/charge corrections for a Dehn
filling.}
\end{center}
\end{figure}

{\it Proof of Theorem \ref{surcomp}.} Again for simplicity, assume
that $L$ is a {\it knot} (one component).  We apply the very same
arguments as for the first claim of Theorem \ref{defscissors}. In
particular, Lemma \ref{extc} applies verbatim. Since $\rho(l)$ is
trivial, for an arbitrary flattening the weight along $l$ (computed in
the flattened $\Ii$-cusps) lies in $\mz\pi\sqrt{-1}$. Hence we can
again use Theorem \ref{theoweight} to deduce the existence of
flattenings with zero weight along $l$. Then we give the same
log-branches on the $3$-simplices of the singular $3$-ball $CD$, in a
pair above and below the disk $D$ (see \cite{N}, p. 454). Note that if
we use $\Dd$ triangulations leading to $\Hh_N$ partition functions,
then both parallel components $L,L'$ that make $\bar \lambda$ survive
in the Hamiltonian subcomplex. If we can deal with $\Hh_N^e$-ones,
only the framing longitude $L'$ survives.  Hence we eventually get a
distinguished flat/charged $\Ii$-triangulation for $(W(l),l^* \cup
\bar{\lambda},\rho)$, or $(W(l),l^* \cup
\lambda',\rho)$ respectively. \hfill$\Box$

\subsection{Manifolds that fiber over $S^1$ - Examples}\label{FIBERED} Lemma
\ref{CYLtensor} (3) gives a practical recipe to compute the QHFT
partition functions of mapping tori. A specific class of
distinguished flat/charged $\Ii$-triangulations of
$(\tilde{W}_{[\psi]},\tilde{L}_{[\psi]},\rho)$ is obtained by
composing one for the trivial mapping cylinder $F \times [-1,1]$, say
$\Tt_{\rm triv}$, with the monodromy action on the $e$-triangulation
$(T,b)$ of $F \times \{1\}$, and then gluing the two boundary
components. The monodromy action can always be decomposed as a
sequence of flip moves: a single flip on the ideal triangulation $T'$
associated to $T$ defines a flip on $(T,b)$ if it is not adjacent to a
marked corner, and it lifts to sequences as in Figure \ref{liftflip}
otherwise. We view these sequences as the result of gluing
tetrahedra. Hence the monodromy action determines a branched
triangulated pseudo-manifold $T_s$. This can be completed with global
charges, and, as for any $\rho$ we are free to choose the cocycle in
$\Tt_{\rm triv}$, we can also complete $T_s$ to a flattened
$\Ii$-triangulation.  Equivalently we can define a sequence
\begin{equation}\label{mon1}
s: (T,b,\Ll) \rightarrow \ldots \rightarrow
\psi(T,b,\Ll)\end{equation} of $e$-triangulations with
($+$)-Log-$\Ii$-parameters compatible with $\rho$. Note that the
edges of the associated pattern $\Tt_{s}$ of flat/charged
$\Ii$-tetrahedra are disjoint from the Hamiltonian link $H$. Since
$\Hh_N(\Tt_{\rm triv})=_N {\rm id}$, we deduce that
$\Hh_N(\tilde{W}_{[\psi]},\tilde{L}_{[\psi]},\rho)=_N
\Hh_N(\Tt_{s})$. 
\begin{figure}[ht]
\begin{center}
\includegraphics[width=6cm]{liftflip.eps}
\caption{\label{liftflip} Lifts to $e$-triangulations (economic ones -
  see Remark \ref{varpar} - at the first row) of flip moves on the
  corresponding ideal triangulations near marked corners. The
  tetrahedron associated to the first flip (first and third ones for
  the second row) degenerates for a sequence $s$ with
  ($-$)-Log-$\Ii$-parameters.}
\end{center}
\end{figure}

Using a similar construction we now prove the
relationship with the quantum hyperbolic invariants $H_N$ of
fibered cusped manifolds \cite{BB2}. Recall that $W\setminus L$ is
homeomorphic to ${\rm Int}(\tilde{W}_{[\psi]})$. Denote by $l$ the
number of components of $L$.
\begin{prop} \label{teofibered} 
If ${\rm Int}(\tilde{W}_{[\psi]})$ supports a (necessarily unique)
complete hyperbolic structure with holonomy $\rho_c$, then
$\Hh_N(\tilde{W}_{[\psi]},\tilde{L}_{[\psi]},\rho_c,0,0)=_N
N^{2l}\ H_N(W\setminus L)$.
\end{prop}

\noindent {\it Proof.} Let $S\setminus V$ be the fiber of $W\setminus
L \rightarrow S^1$, equipped with an ideal triangulation $T'$. First
we show the existence of sequences $s:T'\rightarrow \ldots \rightarrow
\psi(T')$ of flip moves decomposing the monodromy action, such that
the associated pseudo-manifolds $T'_s$ are topological ideal
triangulations of $W\setminus L$, which moreover have maximal
volume. 

The first condition follows from the fact that the monodromy is
homotopically aperiodic (ie. pseudo-Anosov), so that $T'_s$ is
genuinely three-dimensional. When $S\setminus V$ is a once-punctured
torus, the second condition is a consequence of a result of Lackenby
\cite{La}, showing that the {\it monodromy ideal triangulation} of
Floyd and Hatcher \cite{FH} is isotopic to the canonical
Epstein-Penner cellulation. More in general, since no edge of $T'_s$
is homotopically trivial, the results of \cite{F1} imply that we can
straighten the tetrahedra to oriented geodesic ones, possibly with
overlappings, so that the algebraic sum of volumes is Vol($W\setminus
L$). This is known to be maximal \cite{F2}.

As in Section \ref{DEHNF} we can complete $T_s'$ to a flat/charged
$\Ii$-triangulation $\Tt'$. Hence the invariants $H_N(W\setminus L)$
can be computed as trace tensors $\Hh_N(\Tt')$. We note that in the
case when there are several fibrations of $W\setminus L$, or $T_s'$ is
not canonical, the invariance follows from Theorem 6.8 (2) in
\cite{BB2}, which shows that any two flat/charged $\Ii$-triangulations
of $W\setminus L$ with maximal volume are QHG-isomorphic.

Let us denote $\Tt_{s}'$ the result of cutting $\Tt'$ along the
fiber. The two boundary copies are marked pleated hyperbolic surfaces
$(T',b',\mathcal{L}) \rightarrow S\setminus V$ and
$(T',b',\mathcal{L}) \rightarrow \psi(S\setminus V)$, with shear-bend
coordinates (ie. $(-)$-Log-$\Ii$-parameters) $\Ll$ that determine
completely the log-branches of $\Tt_s'$. Recall from Section
\ref{+-versus} the families of flat/charged $\Ii$-triangulations
$\mathcal{F}_t$ and $\mathcal{F}$, and let $C \in \mathcal{F}$, $C_t
\in \mathcal{F}_t$ have boundary structures $(T,b,\mathcal{L})$,
associated to $(T',b',\mathcal{L})$, and $(T,b,\mathcal{L}_t)$,
respectively, at $F\times \{1\}$. We have: $$\begin{array}{lll}
\Hh_N(\Tt') & = & {\rm Tr}(\Hh_N(\Tt_{s}'))\\ & = & N^{-2l}\ {\rm
Tr}\left( \Hh_N(\Tt_{s}') \otimes {\rm id}^{\otimes 2l} \right)\\ & =
& N^{-2l}\ {\rm Tr}\left(\Hh_N(C) \circ (\Hh_N(\Tt_{s}') \otimes {\rm
id}^{\otimes 2l}) \circ \Hh_N(\psi(C))^{-1}\right)\\ & = & N^{-2l}\
\lim_{t \rightarrow {\rm id}} {\rm Tr}\left(\Hh_N(C_t) \circ
(\Hh_N(\Tt_{s,t}')\otimes {\rm id}^{\otimes 2l})\circ
\Hh_N(\psi(C_t))^{-1}\right)\\ & = & N^{-2l}\ \lim_{t \rightarrow {\rm
id}} {\rm Tr}(\Hh_N(C_t * \Tt_{s,t}'*(-\psi(C_t)))).
\end{array}$$
Here we use the invertibility of $\Hh_N(C)$ (Proposition
\ref{tunnel+-}) and the equality $\Hh_N(C) =_N \Hh_N(\psi(C))$. We
define $\Tt_{s,t}'$ as the continuous deformation of $\Tt_{s}'$
obtained from the sequence
$$s_t: (T,b,\Ll_t) \rightarrow \ldots \rightarrow \psi(T,b,\Ll_t)$$
similarly as in (\ref{mon1}) (see also Figure \ref{liftflip}). In the
last equality, $C_t * \Tt_{s,t}'*(-\psi(C_t))$ is for any fixed $t$ a
distinguished flat/charged $\Ii$-triangulation of the mapping cylinder
of $\psi$. Hence $\Hh_N(C_t * \Tt_{s,t}'*(-\psi(C_t)))$ does not
depend on $t$ up to conjugacy, and we conclude with Lemma
\ref{CYLtensor} (3).\hfill$\Box$

\medskip

By following the above computation backwards, we see more in general
that for {\it any} $PSL(2,\mc)$-character $\rho$ that can be realized
by ($-$)-Log-$\Ii$-parameters on some ideal triangulation of the fiber
$S\setminus V$, we have
\begin{equation}\label{mon2}
\Hh_N(\tilde{W}_{[\psi]},\tilde{L}_{[\psi]},\rho,h,k) = N^{2l}\ {\rm
Tr}( \Hh_N(\mathcal{T}_s'))
\end{equation}
where $\mathcal{T}_s'$ is a pattern of flat/charged $\Ii$-tetrahedra
associated to a sequence similar to $(\ref{mon1})$, but with
$e$-triangulations equipped with $(-)$-Log-$\Ii$-parameters compatible
with $\rho$.
\begin{remark}\label{QTeich}{\rm The formula $H_N(W\setminus
L) = {\rm Tr}(\Hh_N(\Tt_{s}'))$ expresses the quantum hyperbolic
invariants of fibered cusped manifolds as amplitudes between two
markings of the fiber, identified with a pleated hyperbolic
surface. For a similar construction based on representations of
quantum Teichmuller spaces, see \cite{BoL}.}
\end{remark}
{\bf Example: the figure-eight knot complement}\label{EXAMPLES}
Here we compute the QHFT partition functions of $(S^3,K_0)$, where
$K_0$ is the $0$-framed figure-eight knot in $S^3$. Recall that
$S^3\setminus K$ is fibered over $S^1$, with fiber the once-punctured
torus $\Sigma_{1,1}$; the $0$-framing of $K$ is induced by the
fibration. For simplicity, below we consider only characters $\rho$ of
{\it injective} representations $\pi_1(S^3\setminus K)\rightarrow
PSL(2,\mc)$. The restriction to $\Sigma_{1,1}$ of such representations
can be realized by $(-)$-Log-$\Ii$-parameters on any ideal
triangulation of $\Sigma_{1,1}$, so that, by (\ref{mon2}), we can
determine the corresponding subspace (still denoted $Def(S^3,K)$) of
the phase space of Definition \ref{defps} by using the monodromy ideal
triangulation.
\smallskip

 The monodromy $\Phi:\Sigma_{1,1}\rightarrow \Sigma_{1,1}$ of
 $S^3\setminus K$ is isotopic to the hyperbolic element
$$\left(  \begin{array}{cc} 2 & 1 \\ 1 & 1  \end{array} \right) =
\left( \begin{array}{cc} 1 & 1 \\ 0 & 1 \end{array}\right) \left(
  \begin{array}{cc} 1 & 0 \\ 1 & 1 \end{array} \right) \quad \in
SL(2,\mz).$$ This description of $[\Phi] \in {\rm Mod}(\Sigma_{1,1})$
can be understood in terms of the Diagram of $PSL(2,\mz)$ (see
eg. \cite{FH}) via the action of $\Phi$ on topological ideal
triangulations of $\Sigma_{1,1}$, which can be represented by two flip
moves. See Figure \ref{Tdyn}, where the left picture is a lift to
$\mr^2 \setminus \mz^2$ of such a triangulation.
\begin{figure}[ht]
\begin{center}
\includegraphics[width=10cm]{Tdyn.eps}
\caption{\label{Tdyn} The composition of flip transformations
  for the monodromy of $S^3\setminus K$.}
\end{center}
\end{figure}

The monodromy ideal triangulation $T$ of $S^3 \setminus K$ is obtained
by realizing each flip move via the gluing of an ideal tetrahedron,
first on a fixed triangulation of $\Sigma_{1,1}$, then on the
resulting one. The remaining four free faces are identified under
$\Phi$. It is not difficult to see that $T$ is isotopic to the
canonical geodesic ideal triangulation of $S^3 \setminus K$ with its
complete hyperbolic structure. The gluing pattern of the tetrahedra in
$T$ is shown in Figure \ref{Teight}.
\begin{figure}[ht]
\begin{center}
\includegraphics[width=8cm]{Teight.eps}
\caption{\label{Teight} The face and edge identifications for the canonical
  geodesic ideal triangulation of $S^3\setminus K$.}
\end{center}
\end{figure}

It is well-known (see \cite{Th}) that the deformation space of smooth
(non necessarily complete) hyperbolic structures on $S^3 \setminus K$
is isomorphic to the algebraic set $Def_{{\rm hyp}}(S^3 \setminus K)
\subset \mathbb{H}^2 \times \mathbb{H}^2$ of points $(w_2,z_0)$ such
that
$$w_2 \in \mathbb{H}^2 \setminus \left\lbrace \frac{1}{2} +
\frac{t}{2}i \mid t \geq \sqrt{15} \right\rbrace,\quad z_0 =
\frac{1}{2} + \left(\frac{1}{4} +
\frac{1}{w_2(w_2-1)}\right)^{\frac{1}{2}}$$ where $w_2$ is the
cross-ratio modulus of the edge $e_2$ in the tetrahedron $\Delta^+$
with positive branching orientation (back edge in the left tetrahedron
of Figure \ref{Teight}), and similarly for $z_0$ in $\Delta^-$. The
space $Def_{{\rm hyp}}(S^3 \setminus K)$ is a subspace of 
\begin{equation}\label{subs}
\mathcal{C}=\{
(w_2,z_0) \in (\mc \setminus \{0,1\})^2 \mid
w_1w_2^2z_0^{-2}z_1^{-1}=1\},
\end{equation} which is isomorphic to the whole set of solutions of the edge
compatibility relations for cross-ratio moduli (see Definition
\ref{non-rough}). By an easy computation we find that the edge
compatibility relations for log-branches and charges are:
\begin{equation}\label{eqfc}
(S)\left\lbrace\begin{array}{l}f_1^-+2f_0^-+2f_0^++f_1^+=
(\arg(w_1)+2\arg(w_2)-\arg(z_1)-2\arg(z_0))-2*_+\\
c_1^-+2c_0^-+2c_0^++c_1^+=0\end{array}\right.
\end{equation}
where $*_+$ is the sign of the imaginary part of the $w_i$, and the
flattenings $f_i^+$ and $f_i^-$ correspond to the cross-ratio moduli
$w_i$ and $z_i$ of $\Delta^+$ and $\Delta^-$, respectively. Hence, by
(\ref{qlb}) we see that the QHFT phase space $Def(S^3,K)$ is the
covering of $\mathcal{C}$ given by
$$Def(S^3,K)=\{((w_2;f_2^+-c_2^+,f_0^+-c_0^+),(z_0;f_0^-+c_0^-,f_1^-+c_1^-))
\in \hat{\mc}\times\hat{\mc} \mid (w_2,z_0) \in \mathcal{C},\ (S){\rm\
is\ satisfied}\ \}.$$ Let us point out some geometrically meaningful
subspaces (compare with \cite{N}, Section 15). The dilation factors of
the standard meridian $m$ and longitude $l$ of $K$ are
$$\mu(m) = z_2w_2^{-1},\quad \mu(l)=z_0^2z_2^{-2}.$$ The complete
hyperbolic structure of $S^3\setminus K$ is obtained by solving
$\mu(m)=1$ and $\mu(l)=1$ for points of $Def_{{\rm hyp}}(S^3 \setminus
K)$, which gives $z_j=(w_j)^* = \exp(i\pi/3)$. We already noticed that
the weights $k$ for this solution take values in $2\mz$ (see \cite{N},
Proposition 5.2). Corresponding flattenings satisfy
\begin{equation}\label{eqcompfc}f_2^--f_2^+ = k(m),\quad2f_0^--2f_2^-=k(l)
\end{equation}
and (\ref{eqfc}) simplifies to
\begin{equation}\label{simpfs}
f_1^-+2f_0^-+2f_0^++f_1^+=0.
\end{equation}
Together with (\ref{eqcompfc}) this gives
\begin{equation}\label{standardf} f_1^+ = \frac{k(l)}
{2}-1-2f_0^+,\quad \left\lbrace \begin{array}{l} f_0^-=k(m)-f_0^+\\
    f_1^- = 2k(m)-\frac{k(l)}{2}+1+2f_0^+\end{array}\right.
\end{equation}
Similarly, integral charges are obtained by letting $c_i^+ = -f_i^+$
on $\Delta^+$ and $c_i^- = +f_i^-$ on $\Delta^-$. It can be checked
that $h=0$. Of course, (\ref{standardf}) is still true for points of
$Def(S^3,K)$ sufficiently near the complete solution, if now $k(m)$
and $k(l)$ are replaced with
$\tilde{k}(m)=k(m)-\log(\mu(m))/\sqrt{-1}\pi$ and
$\tilde{k}(l)=k(l)-\log(\mu(l))/\sqrt{-1}\pi$. Hence, in the vicinity
of the complete structure we have the strata
$$\{((w_2;2f_0^+-\tilde{k}(l),2f_0^+),(z_0;2\tilde{k}(m)-2f_0^+,
2(2\tilde{k}(m)+1+2f_0^+)-\tilde{k}(l)),\ f_0^+ \in \mz\},$$
parametrized by the lifts $k(m)$ and $k(l)$ of $\mu(m)$ and $\mu(l)$.
\smallskip

For points of $\mathcal{C}$ corresponding to hyperbolic structures
whose metric completion is obtained from $S^3 \setminus K$ by $(p,q)$
Dehn filling, we have $p\mu(m) +q\mu(l) = 2\pi \sqrt{-1}$. Hence the
flattenings $f=f''$ of Lemma \ref{deformedlogb} are to satisfy
$$p(f_2^--f_2^+)+q(2f_0^--2f_2^-)=-2$$ Let us fix $r$, $s \in \mz$
such that $ps-qr=1$. Solving simultaneously (\ref{simpfs}) and the
last equation gives
\begin{equation}\label{surgf}
f_1^+ = r-1-2f_0^+,\quad \left\lbrace \begin{array}{l}
    f_0^-=-2s-f_0^+\\ f_1^- = -r+4s+1+2f_0^+\end{array}\right.
\end{equation}
The parity condition $f_0^++f_1^++f_0^-+f_1^- \in 2\mz$ (for the class
$h$ to be $0$) is automatically satisfied.
\smallskip

Now we can compute the QHFT invariants of $(S^3,K_0)$ as functions on
$Def(S^3,K)$. Given a $N$-state $s$ of $T$, put $\alpha=s(\hat{2})$,
$\beta=s(\hat{0})$, $\gamma=s(\hat{3})$ and $\delta=s(\hat{1})$, where
$\hat{i}$ denotes the face of $\Delta ^+$ opposite to the $i$th
vertex. By \ref{mon2}) we have:
$$N^{-2}\Hh_N(S^3,K_0,\rho) = \sum_{\alpha,\beta,\gamma,\delta=0}^{N-1}
  \Rr_N(\Delta^+,b^+,w,f^+,c^+)_{\alpha,\beta}^{\gamma,\delta}\
  \Rr_N(\Delta^-,b^-,z,f^-,c^-)_{\beta,\alpha}^{\delta,\gamma}$$
$$\begin{array}{lll} & = & \sum_{\alpha,\beta,\gamma,\delta=0}^{N-1}
    (w_0'^{-c_1^+}w_1'^{c_0^+})^{\frac{N-1}{2}} \
    \frac{g(w_0')}{g(1)}\ \zeta^{\gamma\beta+(m+1)\gamma^2}\
    \omega(w_0',w_1'^{-1}\vert \alpha-\gamma) \ \delta(\alpha+\beta
    -\delta)\\ & & \hspace{4cm} \times
    (z_0'^{-c_1^-}z_1'^{c_0^-})^{\frac{N-1}{2}}\
    \frac{[z_0']g(1)}{g(z_0')}\ \zeta^{-\alpha\delta-(m+1)\delta^2}\
    \frac{\delta(\alpha+\beta
    -\gamma)}{\omega(z_0'/\zeta,z_1'^{-1}\vert \beta-\delta)} \\ \\ &
    = &
    (w_0'^{-c_1^+}w_1'^{c_0^+}z_0'^{-c_1^-}z_1'^{c_0^-})^{\frac{N-1}{2}}
    \frac{[z_0']g(w_0')}{g(z_0')} \sum_{\alpha,\beta=0}^{N-1}
    \zeta^{\beta^2-\alpha^2}\ \omega(w_0',w_1'^{-1}\vert N-\beta)\
    \omega(z_0'/\zeta,z_1'^{-1}\vert N-\alpha)^{-1}.
\end{array}$$
By the proof of Proposition 8.6 in \cite{BB2} we have
$$\frac{[z_0']g(w_0')}{g(z_0')} =
\frac{g((z_0')^*)^*g(w_0')}{|g(1)|^2}$$
and
$$\omega(z_0'/\zeta,z_1'^{-1}\vert N-\alpha)^{-1} =
\omega((z_0')^*,(z_1'^{-1})^*\vert \alpha)^*$$
where $z^*$ is the complex conjugate of $z$. Thus, setting
$$S(w_0',w_1') = \sum_{\beta=0}^{N-1} \zeta^{\beta^2}
\omega(w_0',w_1'^{-1}\vert \beta) = 1+\ \sum_{\beta=1}^{N-1}
\zeta^{\beta^2} \prod_{k=1}^{\beta} \frac{w_1'^{-1}}{1-w_0'\zeta^k}$$
we get
$$\Hh_N(S^3,K_0,\rho) =
N^{2}(w_0'^{-c_1^+}w_1'^{c_0^+}z_0'^{-c_1^-}z_1'^{c_0^-})^{\frac{N-1}{2}}
\frac{g((z_0')^*)^*g(w_0')}{|g(1)|^2}S(w_0',w_1')
\ S((z_0')^*,(z_1')^*)^*.$$ Using (\ref{standardf}) with
$k(m)=k(l)=f_0^+=0$ and the global charge with $c_0^+=c_0^-=0$, we see
that for the complete hyperbolic structure $\rho_{\rm comp}$ on
$S^3\setminus K$ we have
$$z_0' = (w_0')^*= \exp(i\pi/3N),\quad z_1' = (w_1')^*=
\exp(-5i\pi/3N).$$ Hence $$\Hh_N(S^3,K_0,\rho_{\rm comp})
=N^{2}\frac{|g(e^{i\pi/3N})|^2}{|g(1)|^2} \ |
S(e^{i\pi/3N},e^{-5i\pi/3N})|^2.$$ Let us finally consider hyperbolic
$(p,q)$ Dehn filling of $S^3\setminus K$. Denote $S^3(K_{(p,q)})$ the
surgered manifold, $L$ the core of the surgery, and $\rho_{(p,q)}$ its
hyperbolic holonomy. Because of (\ref{surgf}) the difference $f''-f$
is given on the edges $e_0$, $e_1$ and $e_2$ of $\Delta^+$
(resp. $\Delta^-$) by $0$, $r$ and $-r$ (resp. $-2s$, $4s$ and
$0$). Put $N=2m+1$. From Theorem \ref{defscissors} we deduce
$$\begin{array}{l}\Hh_N(S^3(K_{(p,q)}),L,\rho_{(p,q)}) = N^{2}
(w_0'^{-c_1^+}w_1'^{c_0^+}z_0'^{-c_1^-}z_1'^{c_0^-})^{\frac{N-1}{2}}\
\frac{g((z_0')^*)^*g(w_0')}{|g(1)|^2}\\
\hspace{0.8cm}\sum_{\alpha,\beta=0}^{N-1}
  \zeta^{\beta^2-\alpha^2}\omega(w_0',w_1'^{-1}\vert N-\beta)\
  \omega((z_0')^*,(z_1'^{-1})^*\vert \alpha)^*\
  \zeta^{r(N-\beta)(m+1)}\prod_{j=1}^{N-2s}
  \frac{(z_1'^{-1})^*\zeta^{-4s\alpha(m+1)}}{1-(z_0')^*\zeta^{j+\alpha}}\
  .
\end{array}$$
\begin{remark}{\rm Recall the space $\mathcal{C}$ in (\ref{subs}). 
As already mentionned after Definition \ref{non-rough}, we have a
holonomy map $hol:\mathcal{C}\rightarrow X$ to the character variety
$X=X(\pi_1(S^3\setminus K))$. (See \cite{Ril} or \cite{GAM} for a
complete description of the latter). The map $hol$ is generically
$2:1$, and is onto the geometric component of $X$ \cite{Cha}. We can
express the above partition functions in terms of standard generators
of $X$ by the following observation. Considering $S^3\setminus K$ as
the mapping torus of the monodromy $\Phi$, the edges $e_0$, $e_1$ of
$\Delta^+$ are identified with a longitude $l$ and meridian $m$ of the
punctured torus $\Sigma_{1,1}$, and in $\Delta^-$ we have ($e_0'$ is
opposite to $e_0$):
$$e_0=\Phi(l.m),\quad e_1=\Phi(l),\quad e_0'= \Phi(m).$$ As above,
assume that $\rho$ has non trivial holonomy at $m$, $l$ and
$l.m$. Take a flat/charged $\Ii$-triangulation of $(S^3,K_0,\rho)$ as
in Section \ref{FIBERED}, with $PSL(2,\mc)$-valued cocycle $z$. Denote
$z_l$ the value at $l$, and so on. Note that $z_{\Phi(m)}=Az_mA^{-1}$,
where $A=z(S^1)$, the cocycle value on the standard meridian of the
knot $K$. Then the cross ratio moduli of $\Delta^+$ and $\Delta^-$ are
given by
$$w_0=[0:z_l(0):z_lz_m(0):z_{\Phi(m)}(0)],\quad z_0=
[0:z_{\Phi(l)}z_{\Phi(m)}(0):z_{\Phi(l)}z_{\Phi(m)}z_{\Phi(l)}(0):z_mz_l(0)].$$
(We use the branching to remove the twofold ambiguity of $hol$, as it
allows to specify an equivariant association of a fixed point for each
peripheral subgroup of $\rho(\pi_1(S^3\setminus K))$.)}
\end{remark}
\begin{remark}
{\rm Formulas for Cheeger-Chern-Simons invariants
$\Hh_1(S^3,K_0,\rho)$ of arbitrary $PSL(2,\mc)$-characters of
$S^3\setminus K$ come exactly in the same way (See Remark \ref{OKCCS}
and Remark \ref{scissors}). In the peculiar situation of the complete
hyperbolic structure and its hyperbolic Dehn fillings, they coincide
with those of \cite{N}, Section 15.}
\end{remark}
